\section{Declarative Evidence Algebra}
\label{sec:algebra}

% EA-P1 — Design goals
Evidence acquisition in a multi-hop system is usefully viewed as an execution
problem: a complex question induces a set of evidence conditions that must be
satisfied, and the system must acquire the evidence that satisfies them. Our
algebra is designed around three principles. First, programs are
\emph{declarative}: a program expresses what evidence goals must be met rather
than which retrieval procedure to run. Second, operators are \emph{typed and
composable}: each operator has a typed input/output contract, so programs are
built by composing operators whose evidence types are compatible. Third, the
properties an optimizer needs---expected cost, expected yield, dependency, and
utility---remain \emph{explicit} as typed fields rather than hidden inside
natural-language prompts. The central object is the
\emph{evidence requirement} (\code{EvidenceRequirement}, models.py), a
first-class typed need with a status, an expected evidence type, an importance,
the variables it must bind, and its dependencies over other requirements. Its
status is one of \code{unresolved} $<$ \code{partial} $<$ \code{satisfied}
(\code{RequirementStatus}), forming a monotone lattice over which the execution
semantics advance.

% EA-P2 — Requirement and acquisition operators
We distinguish the declarative condition to be satisfied from the mechanism
used to acquire evidence for it. \textsc{Require} introduces an unresolved
evidence condition; \textsc{Search} and \textsc{Expand} acquire candidate
evidence, either from a source index or from an already-bound value. This
distinction is enforced in the state-transition semantics of the algebra: the
operator \code{apply\_observation(state, observation)} advances each affected
requirement one legal step along the lattice \code{unresolved} $\rightarrow$
\code{partial} $\rightarrow$ \code{satisfied}, and never regresses (a satisfied
requirement is terminal). An observation is a per-requirement coverage probe
(how many of the requirement's variables are bound, whether any evidence has
been seen, and which source ids produced it). Because a requirement encodes
\emph{what must hold}, not \emph{which action to execute}, the same condition
remains well-defined no matter which physical acquisition strategy is later
selected.

% EA-P3 — Binding and composition operators
Two properties make downstream acquisition conditional on observed values
rather than on newly generated subquestions. First, \emph{binding}: evidence is
converted into reusable typed variables that later operators can consume
(\code{BindingRow}, \code{AdaptiveBindingBeam}). Second, \emph{composition}:
operators compose evidence through shared bindings or explicit join
constraints. In the compiled plan, join edges connect the slots that must share
a bound variable; this is what turns a sequence of independent retrieval calls
into a dependent evidence chain. Filtering applies typed predicates to reduce
candidate sets before materialization. The distinction is not cosmetic: because
the dependency is over the variable values actually observed at runtime, a
later acquisition step is conditioned on a real binding (e.g., a bridge entity
that was actually retrieved) rather than on a static subquestion.

% EA-P4 — Verification/materialization/packing operators
A common source of lost information in RAG is the conflation of
\emph{retrieved}, \emph{materialized}, and \emph{packed} evidence. We separate
them deliberately. \textsc{Verify} checks whether candidate evidence supports a
requirement (in our implementation, a sufficiency verdict recorded on the
slot's execution trace); \textsc{Materialize} commits selected evidence to the
persistent evidence state as provenance-clean records; and \textsc{Pack}
constructs the bounded context delivered to the answer model. Evidence can
therefore be retrieved yet rejected before materialization, or materialized yet
omitted from the final context when the reader budget is tighter than the
evidence budget. This separation is what allows a system to report honestly
\emph{how much} of the retrieved evidence survives into the answer context, and
why a gap between retrieval recall and answer quality can be attributed to a
specific stage.

% EA-P5 — Type semantics
The type system exposes five evidence types: passage, entity, relation,
table-row, and structured-record (\code{EvidenceType}). Operator signatures
make compatibility explicit: an operator that consumes a \code{relation}
evidence item is not applied to a bare \code{passage}, and the optimizer can
reason about which physical implementations produce which types. We
additionally establish type-preservation conditions for well-formed programs,
so a legal composition is guaranteed to produce evidence of the declared output
type. We note the current end-to-end backend materializes text passages; the
table-row and structured-record types are the interface over which the
heterogeneous extension (Section~\ref{sec:heterogeneous}) is defined.

% EA-P6 — Rewrite rules
Rewrites fall into three classes, ordered by the strength of their guarantee.
\emph{Exact} rewrites preserve operator semantics exactly. \emph{Evidence-preserving}
rewrites may alter intermediate candidate sets while still guaranteeing that the
requirement is satisfied under stated conditions (e.g., an entity-conditioned
retrieval that returns a superset of the answer-bearing evidence). \emph{Approximate}
rewrites trade recall or verification risk for cost; they expose an estimated
loss term to the optimizer so that a cheaper approximate plan is only selected
when its estimated utility loss is justified. The logical-to-physical mapping
of the next section is precisely an evidence-preserving rewrite family: the
logical program fixes what dependencies must be satisfied, while the physical
implementation (sparse, dense, hybrid, or conditioned access) is left open.

% EA-P7 — End-to-end walkthrough
Figure~\ref{fig:walkthrough} traces a single query through the algebra. The
question is compiled into a typed evidence program (a set of requirements with
a dependency graph); each requirement is then realizable by multiple physical
access paths. The same logical program can therefore be executed by a uniform
static plan, by a cost-only plan, or by the requirement-aware optimized plan of
Section~\ref{sec:optimizer}, under identical retrieval budgets. The logical
layer fixes \emph{what} must be satisfied; the physical layer decides
\emph{how}---and the two are audited separately. This decoupling is the
foundation of the matched-budget evaluation in Section~\ref{sec:experiments}.

\subsection{Heterogeneous Evidence}
\label{sec:heterogeneous}

The evidence algebra types Passage, TableRow, Entity, and Relation are each
backed by a physical implementation: dense/sparse retrieval for passages,
conditioned table access for table rows, entity linking for entities, and KB
joins for relations. The typed algebra ensures that a composition requiring a
\code{TableRow} cannot be satisfied with a bare \code{Passage}, and the
physical planner selects the implementation that produces evidence of the
declared type. End-to-end evidence through table rows is real (Section~\ref{sec:results}),
where FEVEROUS table evidence is converted into typed \code{table\_row}
passages and runs through the same materializer without core-architecture
changes.

\begin{figure}[t]
\centering
\begin{tabular}{@{}l@{}}
\begin{tabular}{|l|l|l|}
\hline
\textbf{Requirement} & \textbf{Status} & \textbf{Type} \\
\hline
S1: \texttt{FindEntity(bridge)} & resolved$\to$partial & passage \\
S2: \texttt{FindEntity(answer)} & unresolved & passage \\
\hline
\end{tabular}
\\[4pt]
$\downarrow$ \texttt{apply\_observation(S1, \{covered\_vars:1, evidence:True\})}$\to$ partial \\
$\downarrow$ \texttt{apply\_observation(S2, \{covered\_vars:1, evidence:True\})}$\to$ satisfied
\end{tabular}
\caption{Walkthrough: a two-hop query is compiled into two typed evidence
requirements with a dependency ($\text{S2}$ depends on $\text{S1}$'s bound
entity). Each observation advances requirements along the lattice;
the logical program is the same regardless of which physical plan executes it.}
\label{fig:walkthrough}
\end{figure}
